\subsection{A.3. Lifting the general adic Cohen algebra}
\label{subsec:A3}

The draft's Definition 4.7 on cote 22/013 is not the later, narrower
complete-local definition in EGA $0_{\rm IV}$ 19.8.1. In particular we
do not assume locality or pseudo finite type. We prove the square-zero
existence assertion in this broader noetherian adic situation.
The scan explicitly restricts Definition 4.7 to noetherian rings.
Its condition P refers back to a preceding section not retained in
this dossier. We therefore state the module condition explicitly in
both standard forms: continuous lifting against discrete modules,
and the stronger existence of a basis of open submodules with
projective quotients. We prove preservation of each form, rather
than silently replacing P by flatness.
Cotangent complexes use cohomological degrees $\leq0$;
$\operatorname{Ext}^i_B(L,N)=\operatorname{Hom}_{D(B)}(L,N[i])$.

\begin{lemma}[Proved: the cotangent complex on each adic quotient]
\label{SA23-A-14}
Let $R\to B$ be a map of noetherian rings, give $R$ the discrete
topology and $B$ the $J$-adic topology, and suppose it is formally
smooth for these topologies. For every $n\geq1$, put $B_n=B/J^n$.
Then
\[
 L_{B/R}\otimes_B^{\mathbf L}B_n
 \simeq(\Omega_{B/R}\otimes_B B_n)[0],
 \qquad \Omega_{B/R}\otimes_B B_n\text{ is projective over }B_n.
\]
No finite generation of this projective module is asserted.
\end{lemma}
\begin{proof}
Set $P_n=\Omega_{B/R}\otimes_B B_n$. Given a surjection $M\to N$
of $B_n$-modules, an $R$-derivation $d:B\to N$ kills $J^{n+1}$:
Leibniz's formula expresses its value on a product of $n+1$ elements
of $J$ as a sum in $J^nN=0$. The map
\[
 B\longrightarrow B_n\ltimes N,\qquad b\longmapsto(\bar b,d(b))
\]
is therefore continuous when the square-zero extension on the right
is discrete. Lift it through $B_n\ltimes M\to B_n\ltimes N$ by formal
smoothness. Its first coordinate is necessarily $b\mapsto\bar b$,
so its second coordinate is a lifted $R$-derivation. The representing
property of $P_n$ proves projectivity.

We calculate the derived fibre at every prime $\mathfrak q\supset J$
of $B$, with $\mathfrak p=\mathfrak q\cap R$ and
$\ell=\kappa(\mathfrak q)$. The local map
$R_{\mathfrak p}\to B_{\mathfrak q}$ is formally smooth for the
maximal-ideal topology of the target. Indeed a continuous map from
$B_{\mathfrak q}$ to a discrete square-zero quotient restricts to a
$J$-adically continuous map from $B$. A lift on $B$ sends every
element outside $\mathfrak q$ to a unit, because its reduction is a
unit and the kernel is nilpotent. It extends to $B_{\mathfrak q}$.
Continuity follows by squaring an ideal killed by the original map.

Let $\widehat R_{\mathfrak p}$ and $\widehat B_{\mathfrak q}$ be the
maximal-adic completions. Completion preserves this formal
smoothness, and the resulting map of complete local rings is regular
\cite[Proposition 15.50.2 and its cited completion lemma]{rootStacksAndre}.
For a noetherian local ring $(T,\mathfrak m)$, flat base change gives
\[
 L_{\widehat T/T}\otimes_{\widehat T}^{\mathbf L}T/\mathfrak m
 \simeq L_{(T/\mathfrak m)/(T/\mathfrak m)}=0,
\]
because $\widehat T\otimes_T T/\mathfrak m=T/\mathfrak m$ and
$T\to\widehat T$ is flat
\cite[Lemmas 10.97.2 and 92.6.2]{rootStacksNoetherianCompletion,rootStacksBaseChange}.
Apply this to $R_{\mathfrak p}$ and
$B_{\mathfrak q}$ and then use transitivity and localization:
\[
 L_{B/R}\otimes_B^{\mathbf L}\ell
 \simeq
 L_{\widehat B_{\mathfrak q}/\widehat R_{\mathfrak p}}
       \otimes_{\widehat B_{\mathfrak q}}^{\mathbf L}\ell .
\]
The cotangent complex of a regular map of noetherian rings is a flat
module in degree zero; an exact opened formulation is
\cite[Theorem 2.10]{rootBlickleFink2025}. Hence this fibre is
concentrated in degree zero.

Put $E_n=L_{B/R}\otimes_B^{\mathbf L}B_n$. Its zeroth cohomology
is $P_n$. Projectivity supplies a map $P_n[0]\to E_n$ inducing the
identity on $H^0$: represent $E_n$ by a nonpositive complex and lift
$P_n$ through its degree-zero cycles onto its zeroth cohomology.
Let $C_n$ be its cone. At every prime of $B_n$ this map becomes an
isomorphism by the preceding fibre calculation, so all derived
residue-field fibres of $C_n$ vanish. Neeman's Lemma 2.12
\cite[p.~527]{rootNeeman1992} implies $C_n=0$. More explicitly, the objects
$X$ with $C_n\otimes^{\mathbf L}_{B_n}X=0$ form a localizing
subcategory containing all residue fields. His Theorem 2.8 says these
fields generate $D(B_n)$ for a commutative noetherian ring, so the
subcategory contains $B_n$. Thus $E_n\simeq P_n[0]$. No
finite-dimensional approximation has been substituted for $B$.
\end{proof}

\begin{lemma}[Proved: vanishing with complete finite coefficients]
\label{SA23-A-15}
Under Lemma~\ref{SA23-A-14}, assume $B$ is separated and complete
for its $J$-adic topology. For every finite $B$-module $N$,
\[
 \operatorname{Ext}^i_B(L_{B/R},N)=0\qquad(i>0).
\]
\end{lemma}
\begin{proof}
A finite module over a noetherian complete adic ring is complete:
$N=\varprojlim_nN/J^nN$. The transition maps are surjective, so this
limit also computes the derived limit. Derived Hom in its second
variable preserves derived limits. Lemma~\ref{SA23-A-14} and
derived restriction--extension adjunction give
\[
\begin{split}
 \mathbf R\operatorname{Hom}_B(L_{B/R},N)
 &\simeq \mathbf R\varprojlim_n
       \mathbf R\operatorname{Hom}_B(L_{B/R},N/J^nN)\\
 &\simeq \mathbf R\varprojlim_n
       \operatorname{Hom}_{B_n}(P_n,N/J^nN)[0].
\end{split}
\]
The Hom transition maps are surjective. An element at level $n$
is a derivation $B\to N/J^nN$; view this target as a
$B_{n+1}$-module and lift the corresponding map from the projective
module $P_{n+1}$ through $N/J^{n+1}N\to N/J^nN$.
This countable tower is Mittag--Leffler. Its derived limit is
concentrated in degree zero, proving the assertion.
\end{proof}

\begin{lemma}[Proved: preservation of topological formal projectivity]
\label{SA23-A-18}
Let $A$ be noetherian, $I^2=0$, and $A_0=A/I$. Let $B$ be a
noetherian flat $A$-algebra with an ideal $Q$, and give
$B_0=B/IB$ the quotient topology. If $B_0$ is formally projective
as a topological $A_0$-module, then $B$ is formally projective as
a topological $A$-module. Here both base rings are discrete, and
formal projectivity means that every continuous linear map to a
discrete module lifts through any surjective linear map of discrete
modules, as in EGA $0_{\rm IV}$, Definition 19.2.1
\cite{algEGAIV1}.
\end{lemma}
\begin{proof}
For a linearly topologized $R$-module $V$ with a decreasing basis
$(V_n)$ of open submodules, put
\[
 T^i_V(M)=\varinjlim_n\operatorname{Ext}^i_R(V/V_n,M).
\]
In particular $T^0_V(M)=\operatorname{Hom}_{R,\mathrm{cont}}(V,M)$
for $M$ discrete. The module $V$ is formally projective if and
only if $T^1_V(M)=0$ for every $M$. In one direction the long
exact sequence of Ext, followed by the exact filtered colimit,
gives the continuous lifting property. Conversely embed $M$ in
an injective $R$-module $E$. Surjectivity of
$T^0_V(E)\to T^0_V(E/M)$ and
$\operatorname{Ext}^1_R(V/V_n,E)=0$ imply $T^1_V(M)=0$.

For an $A_0$-module $M$, the low-degree change-of-rings exact
sequence, natural in $B/Q^n$, contains
\[
\begin{split}
0\longrightarrow
 \operatorname{Ext}^1_{A_0}(B/(IB+Q^n),M)
 &\longrightarrow\operatorname{Ext}^1_A(B/Q^n,M)\\
 &\longrightarrow
 \operatorname{Hom}_{A_0}
 (\operatorname{Tor}_1^A(B/Q^n,A_0),M).
\end{split}
\]
One obtains this sequence by applying derived Hom to the
degree-zero truncation of
$(B/Q^n)\otimes_A^{\mathbf L}A_0$; its degree-minus-one
cohomology is the displayed Tor module. Flatness of $B$ gives
\[
 \operatorname{Tor}_1^A(B/Q^n,A_0)
   =(IB\cap Q^n)/(IQ^n).
\]
Artin--Rees \cite[Lemma 10.51.2]{rootStacksArtinRees} for the
finite $B$-submodule $IB\subset B$ gives a
constant $c$ such that $IB\cap Q^{n+c}\subset Q^nIB=IQ^n$.
Consequently the transition from index $n+c$ to $n$ in this Tor
system is zero. The filtered colimit of the Hom terms is therefore
zero. The colimit of the left-hand Ext terms is zero by the
formal projectivity of $B_0$. Exactness proves
$T^1_B(M)=0$ whenever $IM=0$.

For an arbitrary $A$-module $M$, both $IM$ and $M/IM$ are
$A_0$-modules. The long exact sequences of Ext for
$0\to IM\to M\to M/IM\to0$, followed by filtered colimit,
give $T^1_B(M)=0$. The criterion just proved gives the desired
formal projectivity. No finite generation over $A$ has been used.
\end{proof}

\begin{lemma}[Proved: the stronger projective-quotient condition]
\label{SA23-A-19}
In Lemma~\ref{SA23-A-18}, suppose in addition that $B$ is complete
and separated for its $Q$-adic topology, and that $B_0$ has a
basis of open $A_0$-submodules with projective quotients. Then
$B$ has a basis of open $A$-submodules with projective quotients.
Thus strict formal projectivity in EGA $0_{\rm IV}$, 19.2.3,
is preserved too.
\end{lemma}
\begin{proof}
Write $J=QB_0$. Put $V_0=B_0$ and choose successively open submodules
$V_1\supset V_2\supset\cdots$ of $B_0$, with
$V_n\subset J^nB_0$ and each $B_0/V_n$ projective over $A_0$.
The surjections $B_0/V_{n+1}\to B_0/V_n$ split. Put
$D_{0,n}=\ker(B_0/V_{n+1}\to B_0/V_n)$ for $n\geq0$.
These modules are projective, so completeness identifies $B_0$, topologically
and $A_0$-linearly, with a product
$\prod_{n\geq0}D_{0,n}$ of projective discrete modules, whose
tail starting at index $n$ maps exactly onto $V_n\subset J^nB_0$.
This follows by
choosing a splitting at each step and writing the inverse limit
of the resulting finite products as the full product.

Every projective $A_0$-module $D_0$ lifts to a projective
$A$-module $D$. Indeed, express $D_0$ as the image of an
idempotent $e_0$ on a possibly infinite free module. Lift its
column-finite matrix to an endomorphism $u$ of the corresponding
free $A$-module, put $\delta=u^2-u$, and set
$e=u+(1-2u)\delta$. The endomorphisms $u$ and $\delta$ commute,
$\delta^2=0$, and direct expansion gives $e^2=e$. Its image
reduces to $D_0$ and is projective. This proof applies to every
cardinality of the free basis.

Lift each $D_{0,n}$ to $D_n$. Projectivity lifts its specified
map into $J^nB_0$ through the surjection $Q^nB\to J^nB_0$,
giving $f_n:D_n\to Q^nB$. Completeness permits the continuous
$A$-linear sum
\[
 h:P=\prod_{n\geq0}D_n\longrightarrow B,
 \qquad h((x_n))=\sum_{n\geq0}f_n(x_n).
\]
Because $I$ is finitely generated,
$IP=\prod_n ID_n$; hence $P/IP=\prod_nD_{0,n}$ and the
reduction of $h$ is the chosen isomorphism to $B_0$. The cokernel
of $h$ equals its product with $I$, and therefore is zero since
$I^2=0$. Let $T=\ker h$. Flatness of $B$ makes
$T/IT\to P/IP$ injective. The reduction of $h$ is an
isomorphism, so $T/IT=0$, and again nilpotence gives $T=0$.
Thus $h$ is an algebraic isomorphism.

We verify continuity of the inverse rather than assume an open
mapping theorem. Let $P_{\geq n}$ be a tail of the product. Choose
$u$ such that $J^uB_0$ lies in the reduction of
$h(P_{\geq n})$. Artin--Rees gives $r$ such that
$IB\cap Q^r\subset IQ^u$. Choose $m\geq\max(n,r)$, so that
$h(P_{\geq m})\subset Q^r$, and then choose $s\geq r$ such
that $J^sB_0$ lies in the reduction of $h(P_{\geq m})$.
For $b\in Q^s$ choose $x\in P_{\geq m}$ with
$b-h(x)\in IB$. This difference also lies in $Q^r$, and hence
in $IQ^u$. Every $z\in Q^u$ has the same reduction as
$h(y)$ for some $y\in P_{\geq n}$; since $I^2=0$, this implies
$Iz\subset h(IP_{\geq n})$. Consequently
$b\in h(P_{\geq n})$. We have proved
$Q^s\subset h(P_{\geq n})$, which proves the inverse is
continuous. The images of the tails form the required open
submodules, and their quotients are the finite products of the
projective modules $D_n$.
\end{proof}

\begin{theorem}[Proved: general noetherian adic Cohen lifting]
\label{SA23-A-16}
Let $A\twoheadrightarrow A_0$ be a surjection of noetherian rings
with kernel $I$, where $I^2=0$. Let $A_0\to B_0$ be a noetherian
algebra, separated and complete for a $J$-adic topology, formally
smooth over the discrete ring $A_0$. Then there is a noetherian
$A$-algebra $B$, flat over $A$, with an identification
$B/IB\simeq B_0$. For the inverse image $\widetilde J$ of $J$,
$B$ is separated, complete and formally smooth over $A$.
If $B_0$ is in addition formally projective as a topological
$A_0$-module, then $B$ has the same property over $A$.
The same assertion holds for strict formal projectivity.
Thus the topological module condition in the Cohen-algebra
question is preserved as well.

The lift with its specified reduction is unique up to an
$A$-isomorphism inducing the identity on $B_0$. Its automorphisms
inducing that identity are the additive group
\[
 \operatorname{Der}_{A_0}(B_0,I\otimes_{A_0}B_0).
\]
Uniqueness of an isomorphism is not asserted.
\end{theorem}
\begin{proof}
We first justify $A_0$-flatness of $B_0$ if it has not already been
included in the definition. Completeness implies
$J\subset\operatorname{rad}(B_0)$: $1-jb$ has its geometric-series
inverse for $j\in J$, $b\in B_0$. At every maximal ideal the
localization argument of Lemma~\ref{SA23-A-14} gives a formally
smooth local homomorphism. The local flatness theorem
\cite[Proposition 15.41.5]{rootStacksRegular} makes it flat.
Flatness is local on source and target, so $B_0$ is $A_0$-flat.

Put $N=I\otimes_{A_0}B_0$, a finite $B_0$-module, and specify
$c:I\to N$ by $i\mapsto i\otimes1$. The square-zero extension
theorem for the cotangent complex
\cite[Lemma 92.16.1]{rootStacksDeformations} places the obstruction
to an $A$-algebra extension
\[
 0\longrightarrow N\longrightarrow B\longrightarrow B_0
       \longrightarrow0
\]
inducing $c$ in $\operatorname{Ext}^2_{B_0}(L_{B_0/A_0},N)$.
This group vanishes by Lemma~\ref{SA23-A-15}, so an extension
exists. Its image $IB$ is exactly $N$, since multiplication by
elements of $B$ supplies all $B_0$-multiples of $i\otimes1$.
The map
\[
 I\otimes_A B=I\otimes_{A_0}B_0=N\longrightarrow B
\]
is injective. Thus $\operatorname{Tor}_1^A(A_0,B)=0$.
Together with flatness of $B_0$ over $A_0$, the square-zero
flatness criterion gives $A$-flatness of $B$. Here is the criterion
in this form. For an $A_0$-module $Q$, the low-degree
change-of-rings exact sequence gives
$\operatorname{Tor}_1^A(Q,B)=0$. For an arbitrary $A$-module $Q$,
both $IQ$ and $Q/IQ$ are $A_0$-modules. The Tor sequence of
$0\to IQ\to Q\to Q/IQ\to0$ gives the same vanishing, which is
the module criterion for flatness.

The ring $B$ is noetherian: the image of an ideal in $B_0$ is
finitely generated, and its intersection with $N$ is a submodule
of a finite module over the noetherian ring $B_0$. Lifts of
generators of the image, together with generators of that
intersection, generate the original ideal.

On $N$, the $\widetilde J$-adic module topology is the $J$-adic
one, since $N^2=0$, and is complete. Exactness of completion on
finite modules over the noetherian ring $B$ (Artin--Rees)
\cite[Lemma 10.97.1]{rootStacksNoetherianCompletion}, applied
to the extension sequence, gives
$0\to N\to\widehat B\to B_0\to0$. The comparison from the
original sequence has identity end maps, so
$B\to\widehat B$ is an isomorphism. This proves separation
and completeness.

We verify formal smoothness. Let $M$ be any $B$-module killed by
a power of $\widetilde J$. Both $IM$ and $M/IM$ are
$B_0$-modules killed by a power of $J$. Flat base change, using
$A$-flatness of $B$, identifies
\[
 L_{B/A}\otimes_B^{\mathbf L}B_0\simeq L_{B_0/A_0}.
\]
Lemma~\ref{SA23-A-14} and adjunction imply
$\operatorname{Ext}^1_B(L_{B/A},IM)=0$ and
$\operatorname{Ext}^1_B(L_{B/A},M/IM)=0$.
The long exact sequence gives
$\operatorname{Ext}^1_B(L_{B/A},M)=0$.

This is the required lifting test. For a continuous $A$-map
$u:B\to C/K$ with $C$ discrete and $K^2=0$, pull back
$C\to C/K$ along $u$. The resulting square-zero $A$-algebra
extension of $B$ has kernel $K$, killed as a $B$-module by a power
of $\widetilde J$. Its class in
$\operatorname{Ext}^1_B(L_{B/A},K)$ vanishes, so it splits as an
$A$-algebra and yields a lift to $C$. If $u$ kills
$\widetilde J^s$, the lift kills $\widetilde J^{2s}$ and is
continuous. Hence $B$ is formally smooth for the required topology.

The additional assertion about topological formal projectivity
follows from Lemma~\ref{SA23-A-18}, with $Q=\widetilde J$.
The stronger projective-quotient formulation follows from
Lemma~\ref{SA23-A-19}.
It uses the actual module hypothesis of $B_0$; formal smoothness
and flatness alone do not imply it. For example
$\mathbf Z\to\mathbf Q$ is flat and formally smooth with the
discrete topology, but $\mathbf Q$ is not a projective
$\mathbf Z$-module: every homomorphism from $\mathbf Q$ to
$\mathbf Z$ is zero, so $\mathbf Q$ cannot embed in a free
$\mathbf Z$-module as a projective module would.

Finally the extension theorem classifies isomorphism classes by
$\operatorname{Ext}^1_{B_0}(L_{B_0/A_0},N)$, which vanishes,
and automorphisms by
$\operatorname{Ext}^0_{B_0}(L_{B_0/A_0},N)
=\operatorname{Der}_{A_0}(B_0,N)$.
The isomorphisms preserve inverse images of $J$, hence are
continuous. These derivations are continuous for the $J$-adic
topology on $N$ because $d(J^{r+1})\subset J^rN$.
This also specifies the topological ambiguity.
\end{proof}

\begin{remark}[Known special cases and source distinction]
\label{SA23-A-17}
EGA $0_{\rm IV}$ \cite[19.7.2, printed pp.~108--109]{algEGAIV1} proves the
complete-local case with no separability restriction on residue
extensions; Remark 19.7.3 distinguishes a different, noncomplete
local lifting question. P\'erez Rodr\'iguez
\cite[Corollary 4.3]{rootPerez2008} proves existence for affine
smooth formal schemes, where smoothness includes pseudo finite
type (Definition 1.3). Neither opened result alone gives the
draft's full nonlocal adic assertion. Our proof combines the
cotangent obstruction theorem with a calculation on every adic
quotient. It does not assert a lift for an arbitrary noncomplete
regular local ring with its maximal-adic topology.
\end{remark}

